\scanpage{283}
\mspage{24}

On a le choix \deleted{entre} maintenant entre

a) Une étude plus ou moins systématique des foncteurs entre intervalles. Je préfère attendre pour cela d'avoir besoin de ceci et de cela, et de ne faire que le strict minimum pour les constructions que j'ai en vue.

b) Une description de la composition des chemins.

c) \unsure{Un|Une} formalisme des catégories de chemins \uline{de type non précisé}.

Comme mon objectif est \unsure{avant tout|cependant|maintenant} de décrire de telles catégories $\underline{\mathrm{Ch}}(X)$ \inserted{(sans indice $\tau$ !)}, avec des bonnes propriétés, je commencerai par c, \deleted{et ceci} regarderai ensuite b) comme un test crucial pour l'utilisabilité des notions obtenues, et enfin m'occuperai de a) quand je saurai de quoi j'aurai \unsure{eu|en} besoin.

4. \uline{Catégorie de chemins} $\underline{\mathrm{Ch}}(X)$.

Je veux ``inclure'' \deleted{\unsure{de les|toutes}} les catégories $\underline{\mathrm{Ch}}_\tau(X)$ de chemins de tous les types qu'on veut, dans une $=$ catégorie $\underline{\mathrm{Ch}}(X)$. Rappelons
\[ \underline{\mathrm{Ch}}_\tau(X)=\underline{\mathrm{Hom}}(I_\tau,X), \]
où $I_{\unsure{\tau}}$ est l'\deleted{\unsure{catég}}intervalle (orienté) \unsure{standard|standart} de type $\tau$. \deleted{On va considérer les $\unsure{I}_\tau$ comme des objets}

L'ensemble de tous les types d'intervalles
% archive sheet number: 282
